\documentclass[12pt]{article}
\usepackage[a3paper, margin=1in]{geometry}
\usepackage{amsmath, amssymb, amsthm, ulem}

\begin{document}

\section*{Domácí úkol 4.1}

Vypracoval: Daniel \textit{"Randál"} Ransdorf \hfill Podpis: \rule{4cm}{0.4pt}

\begin{enumerate}
  \item \textbf{Řešení}:
    Rozeberme dva vztahy ze zadání
    \begin{enumerate}
    \item
    První vztah ze zadání určuje jádro matice A. Plyne:
    \[
      \forall s,t \in \mathbb{Z}_7 \,: \quad
      A
      \left(
      s\left(\begin{array}{c}
      3 \\ 0 \\ 1
      \end{array}\right) +
      t\left(\begin{array}{c}
      0 \\ 4 \\ 1
      \end{array}\right)
      \right)
      = 
      \left(\begin{array}{c}
      0 \\ 0
      \end{array}\right)
    \]
    \[
      \Longrightarrow
      A
      \left(\begin{array}{c}
      3s \\ 4t \\ s + t
      \end{array}\right)
      = 
      \left(\begin{array}{c}
      0 \\ 0
      \end{array}\right)
    \]

    Matice A tedy reprezentuje všechny soustavy rovnic s parametrickým řešením výše.
    Jelikož je řešením soustavy dvourozměrný prostor a matice A je třísloupcová,
    musel při řešení Gaussovou eliminací zbýt jen jeden bázový řádek.
    Řekněme, že za pevnou proměnnou byla zvolená proměnná třetího sloupce a první dvě proměnné tedy musely
    být pomocně vyjádřeny parametry. Zapišme, jak vypadá vyřešená Gauss-Jordan eliminovaná matice.
    Úpravami převeďme jen do Gaussovy eliminace,
    tvar poznáme tak, že na pravé straně získáme nuly (pravá strana původně řešené SLR).

    \[
      \left(\begin{array}{ccc|c}
      0 & 0 & 1 & s+t \\
      0 & 0 & 0 & 0 \\ \hline
      1 & 0 & 0 & 3s \\
      0 & 1 & 0 & 4t \\
      \end{array}\right)
      \overset{\text{3.ř} \cdot 2}{\;\sim\;}
      \left(\begin{array}{ccc|c}
      0 & 0 & 1 & s+t \\
      0 & 0 & 0 & 0 \\ \hline
      2 & 0 & 0 & 6s \\
      0 & 1 & 0 & 4t \\
      \end{array}\right)
      \overset{\text{4.ř} \cdot 5}{\;\sim\;}
      \left(\begin{array}{ccc|c}
      0 & 0 & 1 & s+t \\
      0 & 0 & 0 & 0 \\ \hline
      2 & 0 & 0 & 6s \\
      0 & 5 & 0 & 6t \\
      \end{array}\right)
      \overset{\text{1.ř } + \text{3.ř}}{\;\sim\;}
      \left(\begin{array}{ccc|c}
      2 & 0 & 1 & t \\
      0 & 0 & 0 & 0 \\ \hline
      2 & 0 & 0 & 6s \\
      0 & 5 & 0 & 6t \\
      \end{array}\right)
      \overset{\text{1.ř } + \text{4.ř}}{\;\sim\;}
      \left(\begin{array}{ccc|c}
      2 & 5 & 1 & 0 \\
      0 & 0 & 0 & 0 \\ \hline
      2 & 0 & 0 & 6s \\
      0 & 5 & 0 & 6t \\
      \end{array}\right)
    \]

    Gaussova eliminace převádí libovolnou matici A na matici E posloupností elementárních úprav.
    Každá taková úprava je lineární kombinace řádkových vektorů.
    Tím se zachovává prostor, takže výsledné řádky jsou vyjádřitelné pomocí původních a naopak.
    Když jsme zjistili, že jeden výsledek Gaussovy eliminace je zůstatek jediného nenulového řádku,
    znamená to, že i všechny původní řádky museli být jen násobky tohoto řádku,
    jinak by zbylo řádků po eliminaci více. Jediné, co je třeba pohlídat je, že nejsou
    všechny řádky nulové.
    
    Označme si tedy řádkový vektor $\vec{v} = \left(\begin{array}{ccc} 2 & 5 & 1 \end{array}\right)$

    Určeme si první omezující podmínku pro matici A:

    \[
      A \in \left\{
        \left(\begin{array}{c} k_1 \cdot \vec{v} \\ \hline k_2 \cdot \vec{v} \end{array}\right)
        \quad \middle| \quad (k_1, k_2) \in \mathbb{Z}_7^2 \setminus \left\{\left(0,0\right)\right\}
      \right\}
    \]


    \item
    Druhý vztah ze zadání určuje obraz matice A. Plyne:

    \[
      \forall a,b,c \in \mathbb{Z}_7 \,\exists ! t \in \mathbb{Z}_7 \,: \quad A \left( \begin{array}{c}
        a \\ b \\ c
      \end{array} \right) = t \left( \begin{array}{c}
        3 \\ 4
      \end{array} \right) = \left( \begin{array}{c}
        3t \\ 4t
      \end{array} \right)
    \]

    Z bodu 1.a víme, že matice A se musí dát zapsat pomocí řádkových vektorů $\vec{v}$.
    Vyřešme tedy, které $(k_1, k_2) \in \mathbb{Z}_7^2 \setminus \left\{\left(0,0\right)\right\}$
    vyhovují i druhému vztahu ze zadání:


    \begin{align*}
      \forall a,b,c \in \mathbb{Z}_7 \,\exists ! t \in \mathbb{Z}_7 \,: \quad 
      \left(\begin{array}{c} k_1 \cdot \vec{v} \\ \hline k_2 \cdot \vec{v} \end{array}\right)
      \left( \begin{array}{c}
        a \\ b \\ c
      \end{array} \right) &= \left( \begin{array}{c}
        3t \\ 4t
      \end{array} \right) \\
      \Longrightarrow 
      \left(\begin{array}{c} k_1 \cdot \vec{v} \cdot 
        \left( \begin{array}{ccc} a & b & c \end{array} \right) \\ \hline k_2 \cdot \vec{v} \cdot 
        \left( \begin{array}{ccc} a & b & c \end{array} \right) \end{array}\right)
      &= \left( \begin{array}{c}
        3t \\ 4t
      \end{array} \right) \\
      \text{Označme součin } z = \vec{v} \cdot \left( \begin{array}{ccc} a & b & c \end{array} \right) \\
      \Longrightarrow \left(\begin{array}{c} k_1 \cdot z \\ \hline k_2 \cdot z \end{array}\right)
      &= \left( \begin{array}{c}
        3t \\ 4t
      \end{array} \right)
    \end{align*}

    Získáváme:
    \begin{align*}
      \forall z \in \mathbb{Z}_7 \,\exists ! t \in \mathbb{Z}_7 \,: \quad k_1 \cdot z &= 3t \land k_2 \cdot z = 4t \\
      \Longrightarrow 5k_1 \cdot z &= t \land 2k_2 \cdot z = t \\
      \Longrightarrow 5k_1 \cdot z &= 2k_2 \cdot z \\
      \text{Pro } z=0 \text{ jsou zjevně řešením libovolné } k_1, k_2 &\text{ vyšetřeme proto více omezující rovnici, kde } z \neq 0 \\
      \overset{\cdot z^{-1}}{\;\Longrightarrow\;} 5k_1 = 2k_2 \\
      \Longrightarrow k_1 = 6k_2 \\
      \text{Z 1.a víme, že oba parametry $k_1,k_2$ nesmí být zároveň rovny 0. Zvolme tedy parametr } &\alpha \in \mathbb{Z}_7 \setminus \{0\}
      \text{ a položme } k_2 = \alpha: \\
      \Longrightarrow k_1 = 6\alpha
    \end{align*}

    \end{enumerate}

    Poskládejme řešení dohromady:

    \[
      A = \left(\begin{array}{c} k_1 \cdot \vec{v} \\ \hline k_2 \cdot \vec{v} \end{array}\right)
      = \left(\begin{array}{c} 6\alpha \cdot \left(\begin{array}{ccc} 2 & 5 & 1 \end{array}\right)
         \\ \hline \alpha \cdot \left(\begin{array}{ccc} 2 & 5 & 1 \end{array}\right) \end{array}
      \right)
      = \underline{\underline{\alpha \left(\begin{array}{ccc}
          5 & 2 & 6 \\
          2 & 5 & 1
      \end{array}\right)}} \quad , \alpha \in \mathbb{Z}_7 \setminus \{0\}
    \]

    Rozepsaná řešení:

    \[
    A \in \left\{
      \left(\begin{array}{ccc}
          5 & 2 & 6 \\
          2 & 5 & 1
      \end{array}\right),
      \left(\begin{array}{ccc}
          3 & 4 & 5 \\
          4 & 3 & 2
      \end{array}\right),
      \left(\begin{array}{ccc}
          1 & 6 & 4 \\
          6 & 1 & 3
      \end{array}\right),
      \left(\begin{array}{ccc}
          6 & 1 & 3 \\
          1 & 6 & 4
      \end{array}\right),
      \left(\begin{array}{ccc}
          4 & 3 & 2 \\
          3 & 4 & 5
      \end{array}\right),
      \left(\begin{array}{ccc}
          2 & 5 & 1 \\
          5 & 2 & 6
      \end{array}\right)
    \right\}
    \]


\end{enumerate}

\end{document}
